\chapter{Родительская допустимость обменного моста}

Первая проверка построила конечномерный мост $K$, обращающий сбалансированную
пару Тёплица. Теперь проверяется, является ли этот мост уже допустимой
физической одноформой проекта или требует расширения родителя.

\section{Совместимость с обменной вещественной структурой}

Пусть $R$ меняет две ориентированные ветви местами и комплексно сопрягает
коэффициенты. Для
\begin{equation}
 T_\lambda=
 \begin{pmatrix}S_q&\lambda P_q\\0&S_q^*\end{pmatrix}
\end{equation}
непосредственно выполнено
\begin{equation}
 RT_\lambda=T_\lambda^*R.
 \label{eq:v6-bridge-real-exchange}
\end{equation}
Поэтому для внешне градуированного оператора
\begin{equation}
 F_\lambda=
 \begin{pmatrix}0&T_\lambda^*\\T_\lambda&0\end{pmatrix}
\end{equation}
можно выбрать $J$, меняющий градуированные половины и содержащий $R$, так
что
\begin{equation}
 JF_\lambda=F_\lambda J,
 \qquad J\gamma=-\gamma J.
\end{equation}
Базовые знаки KO6 мост не нарушают. Полная Clifford-линейная реализация
остаётся условной на уже выбранное отдельное Clifford-действие, но нового
знакового препятствия не возникает.

\section{Где живёт проектор моста}

В неограниченном расширении Тёплица поляризация Харди даёт
\begin{equation}
 P_q=(I-S_qS_q^*).
\end{equation}
Это элемент компактного идеала расширения. Если гладкая координатная
алгебра содержит конечноматричные операторы, проектор действительно
представим одноформой. На модуле с оператором числа $N_+$ положим
\begin{equation}
 E_{ij}=|e_i\rangle\langle e_j|.
\end{equation}
Тогда
\begin{equation}
 E_{01}[N_+,E_{10}]=E_{00}=P_0.
 \label{eq:v6-p0-as-compact-oneform}
\end{equation}

Литература по спектральным тройкам тёплицева типа подчёркивает, что
компактный идеал является дополнительной частью расширенной алгебры, а не
частью одного символа $C(S^1)$; см. \path{arXiv:0709.4310} и обзор
\path{arXiv:1911.05823}. Поэтому результат зависит от объявления
координатной алгебры:
\begin{enumerate}
 \item в символической алгебре окружности $P_0$ отсутствует;
 \item в гладком компактном идеале $P_0$ является представленной
 одноформой;
 \item в полном расширении Тёплица обе части сосуществуют, но их физический
 статус должен быть задан родительской мерой.
\end{enumerate}

\section{Сопоставление с физическим пакетом \texorpdfstring{$H_{15}$}{H15}}

Замороженный физический модуль одноформ Тома V содержит ровно три
заряженных типа рёбер $u,d,e$ и имеет пространство кратностей
\begin{equation}
 \mathcal Y_\rho=\mathcal E_\rho\otimes\Lambda_{\rm ch},
 \qquad \Lambda_{\rm ch}\simeq\mathbb C^3.
\end{equation}
Обменный мост действует не между этими тремя рёбрами, а между двумя
ориентациями вспомогательного тёплицева цикла. Поэтому
\begin{equation}
 K\notin\mathcal Y_\rho
\end{equation}
для уже замороженного физического исчисления. Его добавление как четвёртого
нейтрального ребра изменило бы модуль одноформ и потребовало бы повторения
условий первого порядка, junk-фактора и нормировки следа.

Это не уничтожает механизм. Существуют две строго различимые возможности:
\begin{description}
 \item[расширение алгебры] компактный идеал Тёплица объявляется частью
 гладкой координатной алгебры;
 \item[суперсвязностное чтение] $K$ является нечётным эндоморфизмом
 соответствия, а не обычной одноформой Стандартной модели.
\end{description}
Обе возможности являются новыми родительскими данными и не могут быть
вставлены молча.

\begin{theorem}[Статус обменного моста]
Мост совместим с базовым обменным $J$ и реализуется представленной
одноформой гладкого компактного идеала Тёплица. Он не принадлежит ранее
зафиксированному физическому бимодулю одноформ $H_{15}$. Следовательно,
механизм проявления скрытого класса математически открыт, но требует
явного выбора между расширением координатной алгебры и нечётным
суперсвязностным эндоморфизмом.
\end{theorem}

\begin{nogocriterion}[Запрет скрытого добавления четвёртого ребра]
Нельзя считать мост уже выведенным из $H_{15}$ только потому, что он имеет
ранг пятнадцать и сохраняет обмен. Его физический статус должен следовать
из нового родительского выбора и пройти повторный дифференциальный аудит.
\end{nogocriterion}

\begin{protocol}
\begin{enumerate}
 \item обменное условие $RT_\lambda=T_\lambda^*R$: \textbf{выполнено};
 \item базовые знаки $J$, $\gamma$: \textbf{совместимы};
 \item $P_0$ как одноформа компактного идеала: \textbf{построен};
 \item $P_0$ в одной символической алгебре $C(S^1)$: \textbf{нет};
 \item мост в физическом модуле $\mathcal Y_\rho$ на $H_{15}$:
 \textbf{нет};
 \item возможность тёплицева расширения: \textbf{открыта};
 \item возможность нечётного эндоморфизма суперсвязности:
 \textbf{открыта};
 \item мост как уже выведенная физическая мода: \textbf{нет};
 \item следующая проверка:
 \textbf{выбор минимального родителя обменного моста}.
\end{enumerate}
\end{protocol}

Машинный сертификат сохранён в
\path{s2t_v6_exchange_bridge_parent_admissibility_gate_results.json}.